\centerline{\bf \Huge Планиметрия}

\problem{\textbf{1}} В треугольнике $A B C$ проведена биссектриса $A M$. Прямая, проходящая через вершину $B$ перпендикулярно $A M$, пересекает сторону $A C$ в точке $N . A B=6 ; B C=5 ; A C=9$.\\
a) Докажите, что биссектриса угла $C$ делит отрезок $M N$ пополам.\\
б) Пусть $P$ - точка пересечения биссектрис треугольника $A B C$. Найдите отношение $A P: P N$.

\problem{\textbf{2}} В равнобедренном треугольнике $A B C$ с углом $120^{\circ}$ при вершине $A$ проведена биссектриса $B D$. В треугольник $A B C$ вписан прямоугольник $D E F H$ так, что сторона $F H$ лежит на отрезке $B C$, а вершина $E$ - на отрезке $A B$.\\
a) Докажите, что $F H=2 D H$.\\
б) Найдите площадь прямоугольника $D E F H$, если $A B=4$.

\problem{\textbf{3}} Из вершин $A$ и $B$ тупоугольного треугольника $A B C$ проведены высоты $B Q$ и $A H$. Известно, что угол $B$ - тупой, $B C: C H=4: 5, B H=B Q$.\\
а) Докажите, что диаметр описанной вокруг треугольника $A B Q$ окружности в $\dfrac{2 \sqrt{6}}{3}$ раз больше $B Q$.\\
б) Найдите площадь четырехугольника $A H B Q$, если площадь треугольника $H Q C$ равна 25 .

\problem{\textbf{4}} На сторонах $A B, B C$ и $A C$ треугольника $A B C$ отмечены точки $C_1, A_1$ и $B_1$ соответственно, причём $A C_1: C_1 B=8: 3$, $B A_1: A_1 C=1: 2, C B_1: B_1 A=3: 1$. Отрезки $B B_1$ и $C C_1$ пересекаются в точке $D$.\\
a) Докажите, что $A D A_1 B_1$ - параллелограмм.\\
б) Найдите $C D$, если отрезки $A D$ и $B C$ перпендикулярны, $A C=28, B C=18$.

\problem{\textbf{5}} Дан параллелограмм $A B C D$ с острым углом $A$. На продолжении стороны $A D$ за точку $D$ взята точка $N$ такая, что $C N=C D$, а на продолжении стороны $C D$ за точку $D$ взята такая точка $M$, что $A D=A M$.\\
a) Докажите, что $B M=B N$.\\
б) Найдите $M N$, если $A C=7, \sin \angle B A D=\frac{7}{25}$.

\problem{\textbf{6}} Дан ромб $A B C D$. Прямая, перпендикулярная стороне $A D$, пересекает его диагональ $A C$ в точке $M$, диагональ $B D$ - в точке $N$, причем $A M: M C=1: 2, B N: N D=1: 3$.\\
a) Докажите, что $\cos \angle B A D=0,2$.\\
б) Найдите площадь ромба, если $M N=5$.

\newpage
\hspace{3mm}

\problem{\textbf{7}} Около остроугольного треугольника $A B C$ описана окружность с центром $O$. На продолжении отрезка $A O$ за точку $O$ отмечена точка $K$ так, что $\angle B A C+\angle A K C=90^{\circ}$.\\
a) Докажите, что четырёхугольник $O B K C$ вписанный.\\
б) Найдите радиус окружности, описанной около четырёхугольника $O B K C$, если $\cos \angle B A C=\dfrac{3}{5}$, а $B C=48$.

\problem{\textbf{8}} В остроугольном треугольнике $A B C, \angle A=60^{\circ}$. Высоты $B N$ и $C M$ треугольника $A B C$ пересекаются в точке $H$. Точка $O$ - центр окружности, описанной около $\triangle A B C$.\\
a) Докажите, что $A H=A O$.\\
б) Найдите площадь $\triangle A H O$, если $B C=6 \sqrt{3}, \angle A B C=45^{\circ}$.

\problem{\textbf{9}} В прямоугольном треугольнике $A B C$ проведена высота $C H$ из вершины прямого угла $C$. В треугольники $A C H$ и $B C H$ вписаны окружности с центрами $O_1$ и $O_2$ соответственно, касающиеся прямой $C H$ в точках $M$ и $N$ соответственно.\\
a) Докажите, что прямые $A O_1$ и $C O_2$ перпендикулярны.\\
б) Найдите площадь четырёхугольника $M O_1 \mathrm{NO}_2$, если $A C=20$ и $B C=15$.

\problem{\textbf{10}} Окружность, вписанная в трапецию $A B C D$, касается ее боковых сторон $A B$ и $C D$ в точках $M$ и $N$ соответственно. Известно, что $A M=8 M B$ и $D N=2 C N$.\\
a) Докажите, что $A D=4 B C$.\\
б) Найдите длину отрезка $M N$, если радиус окружности равен $\sqrt{6}$.

\problem{\textbf{11}} На стороне острого угла с вершиной $A$ отмечена точка $B$. Из точки $B$ на биссектрису и другую сторону угла опущены перпендикуляры $B C$ и $B D$ соответственно.\\
a) Докажите, что $A C^2+C D^2=A D^2+D B^2$.\\
б) Прямые $A C$ и $B D$ пересекаются в точке $T$. Найдите отношение $A T: T C$, если $\cos \angle A B C=\dfrac{3}{8}$.

\problem{\textbf{12}} Биссектриса угла $A$ треугольника $A B C$ пересекает сторону $B C$ в точке $K$, а окружность описанную около треугольника $A B C$, - в точке $M$.\\
a) Докажите, что треугольник $B M C$ равнобедренный.\\
б) Найдите радиус окружности, описанной около треугольника $K M C$, если $A C=6, B C=7, A B=8$.\\

\centerline{\textbf{180 минут}}

